\documentclass[exercice,a4]{thedocument}%
\usepackage{pstricks}%
\usepackage{array}%
\usepackage{scratch3}%
\usepackage{enumitem}%
\usepackage[french]{babel}%
\usepackage{multido}%
\usepackage{pst-plot}%
\startdatakeys
    \source{24GENMATPO3}
    \BO{2026}
    \cycle{4}
    \competencesDuSocle{CA,RA}
    \objectifsApprentissage{C4A1aa,C4A2ad,C4B2cb,C4D2aa,C4D3ac}
    \motsCles{volume,débit,fonction,antécédent,graphique}
\enddatakeys
\begin{document}%
    


\medskip

La piscine du camping \og le Rocher\fg{} dispose d'un bassin circulaire de forme cylindrique de rayon $3,60$~m et de hauteur $1,50$~m. En fin de saison, on utilise une pompe dont le débit est de $14,1$~m$^3$/h pour vider l'eau de la piscine.

\medskip

\begin{enumerate}
\item Montrer que le volume du bassin, arrondi au dixième de m$^3$, est 61,1m$^3$.
\item Le bassin est plein. On met en route la pompe. Au bout de 2 heures, quel volume d'eau en m$^3$ reste-t-il à vider ?
	\end{enumerate}

On considère la fonction $V\::\: t \longmapsto  61,1 - 0,235 t$.

	\begin{enumerate}[resume]
\item 
	\begin{enumerate}
		\item Montrer que l'expression $V(t)$ permet de déterminer le volume d'eau en m$^3$ qu'il reste à vider dans le bassin en fonction de la durée $t$, exprimée en minute, d'utilisation de la pompe.
		\item Calculer le temps nécessaire pour que le volume d'eau restant à vider soit égal à $30$ m$^3$.
		
On donnera une valeur approchée à la minute près.
	\end{enumerate}
	
\item On a tracé ci-dessous une partie de la représentation graphique de la fonction $V$.

\begin{center}
\psset{unit=0.04cm,arrowsize=3pt 3}
\begin{pspicture*}(-20,-20)(330,120)
\psgrid[unit=2cm,subgriddiv=5,gridlabels=0,gridcolor=gray,subgridcolor=lightgray](-1,-1)(7,3)
\psaxes[linewidth=1.25pt,Dx=50,Dy=50,labelFontSize=\scriptstyle]{->}(0,0)(-20,-20)(330,120)
%\multido{\n=0+10}{32}{\psline[linewidth=0.1pt](\n,0)(\n,110)}
%\multido{\n=0+10}{12}{\psline[linewidth=0.1pt](0,\n)(330,\n)}
%\multido{\n=0+50}{6}{\psline[linewidth=0.7pt](\n,0)(\n,110)}
%\multido{\n=0+50}{3}{\psline[linewidth=0.7pt](0,\n)(330,\n)}
\psplot[plotpoints=2000,linewidth=1.25pt,linecolor=blue]{-20}{330}{61.1 0.235 x mul sub}
\uput[u](305,0){$t$}\uput[r](0,105){$V(t)$}
\end{pspicture*}
\end{center}

Répondre aux questions suivantes par une lecture graphique.
	\begin{enumerate}
		\item Déterminer l'antécédent de 40 par la fonction $V$. Interpréter le résultat.
		\item Déterminer le temps nécessaire pour que la pompe vide complètement le bassin.
	\end{enumerate}
\end{enumerate}

\newpage

\begin{center}\textbf{\large ANNEXE}\end{center}

\medskip


\vspace{1cm}

\begin{tabular}{c c c}
\textbf{Question 2. a.}&&\textbf{Question 2. b.}\\%[7pt]
\\
\begin{scratch}[scale=1]
\initmoreblocks{définir \namemoreblocks{pentagone}}
\blockpen{stylo en position d'écriture}
\blockrepeat{répéter \ovalnum{\hspace*{0.7cm}} fois}
{\blockmove{avancer de \ovalnum{longueur} pas}
\blockmove{tourner \turnleft{} de \ovalnum{\hspace*{0.7cm}} degrés}
}
\blockpen{relever le stylo}
\end{scratch}
&&
{\renewcommand{\arraystretch}{1.6}
\begin{tabular}[t]{|m{1.5cm}|m{1.3cm}|m{3.5cm}|}
\multicolumn{3}{c}{~~}\\
\hline
Nom de\newline l'élève&Numéro\newline de la\newline copie d'écran&Nom de la\newline transformation\\ \hline
Camille&&\\ \hline 
Lou &&\\ \hline 
Zoé&&\\ \hline 
\end{tabular}}\\
\end{tabular}

\vspace{1cm}

\textbf{Question 2. c.}

\bigskip

\begin{center}
\begin{tabular}{|l| p{7cm}|}\hline
Instruction&Ordre d'apparition de l'instruction dans le programme de Sofia\\ \hline
\begin{scratch}[scale=0.8]\blockpen{effacer tout} \end{scratch}&\\ \hline
\begin{scratch}[scale=0.8]\blockmove{s'orienter à \ovalnum{90} } \end{scratch}&\\ \hline
\begin{scratch}[scale=0.8]\blockmoreblocks{pentagone} \end{scratch}&\multicolumn{1}{|c|}{6\up{e}}\\ \hline
\begin{scratch}[scale=0.8]\blockinit{Quand \greenflag est cliqué} \end{scratch}&\multicolumn{1}{|c|}{1\up{re}}\\ \hline
\begin{scratch}[scale=0.8]\blockvariable{mettre \selectmenu{longueur} à \ovalnum{60}} \end{scratch}&\\ \hline
\begin{scratch}[scale=0.8]\blockmove{aller à x : \ovalnum{0} y : \ovalnum{0}} \end{scratch}&\\ \hline
\begin{scratch}[scale=0.8]\blockmoreblocks{pentagone} \end{scratch}&\\ \hline
\begin{scratch}[scale=0.8]\blockvariable{mettre \selectmenu{longueur} à \ovaloperator{\ovalvariable{longueur} *\ovalnum{1.5}}} \end{scratch}&\\ \hline
\end{tabular}
\end{center}
\showthesource
\end{document}
